\subsection{Graphs}

\subsubsection*{General Properties}
\begin{itemize}
    \item \textbf{Handshaking Lemma: } The sum of the degrees of all vertices is exactly twice the number of edges ($\sum \deg(v) = 2|E|$).
    \item \textbf{Bipartite Graph: } A graph whose vertices can be divided into two disjoint sets such that every edge connects a vertex in one set to a vertex in the other. A graph is bipartite if and only if it is \textbf{2-colorable} and contains \textbf{no cycles of odd length}.
    \item \textbf{Complete Graph ($K_n$): } A simple graph where every pair of vertices is connected by an edge. It contains exactly $\frac{N(N-1)}{2}$ edges.
    \item \textbf{DAG (Directed Acyclic Graph): } A directed graph with no directed cycles. Every DAG has at least one valid \textit{Topological Ordering}.
    \item \textbf{Planar Graph: } A graph that can be drawn without edges crossing. \textbf{Euler's Formula:} For a planar graph with $V$ vertices, $E$ edges, $F$ faces, and $C$ connected components: $V - E + F = 1 + C$.
\end{itemize}

\subsubsection*{Tree Properties}
\begin{itemize}
    \item \textbf{Definition \& Edges: } A tree is a connected graph with no cycles. A tree with $N$ vertices always has exactly $N-1$ edges.
    \item \textbf{Unique Paths: } There is exactly one simple path between any two vertices in a tree.
    \item \textbf{Cycle Creation: } Adding exactly one new edge to any tree will create exactly one cycle.
    \item \textbf{Bipartite Nature: } Every tree is inherently a bipartite graph (since it has no cycles, it has no odd cycles).
\end{itemize}

\subsubsection*{Important Graph Theorems}
\begin{itemize}
    \item \textbf{Eulerian Cycle \& Path: } 
        \begin{itemize}
            \item An undirected connected graph has an \textbf{Eulerian Cycle} (visits every edge exactly once and returns to start) $\iff$ \textbf{every} vertex has an even degree.
            \item It has an \textbf{Eulerian Path} (visits every edge exactly once but doesn't return) $\iff$ exactly \textbf{zero or two} vertices have an odd degree.
        \end{itemize}
    \item \textbf{Dirac's Theorem (Hamiltonian Cycle): } In a simple graph with $N \ge 3$ vertices, if every vertex has a degree $\deg(v) \ge \frac{N}{2}$, the graph is guaranteed to have a Hamiltonian cycle (visits every vertex exactly once).
    \item \textbf{Ore's Theorem (Hamiltonian Cycle): } In a simple graph with $N \ge 3$ vertices, if $\deg(u) + \deg(v) \ge N$ for every pair of non-adjacent vertices $u$ and $v$, the graph contains a Hamiltonian cycle.
    \item \textbf{Mantel's Theorem: } The maximum number of edges in an $N$-vertex graph that does not contain any triangles (cycles of length 3) is strictly bounded by $\lfloor \frac{N^2}{4} \rfloor$.
    \item \textbf{Kőnig's Theorem: } In any \textbf{bipartite graph}, the size of the Maximum Bipartite Matching is exactly equal to the size of the Minimum Vertex Cover. \textit{(Reduces an NP-Hard problem to a Max Flow / Bipartite Matching problem).}
    \item \textbf{Dilworth's Theorem: } In any DAG, the size of the maximum anti-chain (largest set of vertices where none can reach another) is equal to the minimum number of directed paths needed to cover all vertices.
\end{itemize}